\section{Energieverteilungsnetze}

\subsection{Drehstrom (DS), 3-Phasen-System}



\subsubsection{Spannungen und Ströme in DS}
    \underline{Leiter-Erde-Spannung \(U_{LE}=\SI{230}{\volt}\)}
    \begin{flalign*}
        U_{LE} &= U_Y &\\
        \underline{U}_{L1} &= U_{LE} \angle 0^\circ &\\
        \underline{U}_{L2} &= U_{LE} \angle -120^\circ &\\
        \underline{U}_{L3} &= U_{LE} \angle -240^\circ &
    \end{flalign*}

    \underline{Leiter-Leiter-Spannung \(U_{LL}=\SI{400}{\volt}\)}
    \begin{flalign*}
        U_{LL} &= U_\Delta = U_{LE} \cdot \sqrt{3} &\\
        \underline{U}_{12} &= \underline{U}_{L1} - \underline{U}_{L2} = U_{LL} \angle 30^\circ &\\
        \underline{U}_{23} &= \underline{U}_{L2} - \underline{U}_{L3} = U_{LL} \angle -90^\circ &\\
        \underline{U}_{31} &= \underline{U}_{L3} - \underline{U}_{L1} = U_{LL} \angle -210^\circ &
    \end{flalign*}

    \underline{Ströme in DS (symmetrisch)}
    \begin{flalign*}
        \underline{I}_{Lx} &= \dfrac{\underline{U}_{Lx}}{\underline{Z}_Y} &\\
        \underline{I}_N &= \underline{I}_{L1} + \underline{I}_{L2} + \underline{I}_{L3} = 0 &
    \end{flalign*}

    \begin{table}[H]
        \begin{tabular}{l l l}
            \toprule
            \textbf{Stranggröße} & \textbf{Stern} & \textbf{Dreieck} \\
            \midrule
            Spannung & \(U_{LE} = \frac{U_{LL}}{\sqrt{3}}\) & \(U_{LL}\) \\
            Strom & \(I_L\) & \(I_L = \sqrt{3} \cdot I_{str}\) \\
        \end{tabular}
    \end{table}

\subsubsection{Einphasiges ESB (symmetrisch)}
    Es werden die Leiter-Erde-Spannungen \(U_{LE}\) verwendet. Die Ströme sind die Strangströme \(I_{str}\).
    Im einphasigen ESB ist aufgrund der symmetrischen Belastung (\(I_N = 0\)) der Neutralleiter nicht notwendig.

    Um eine Dreieckschaltung so darzustellen, wird eine Dreieck-Stern-Umwandlung durchgeführt. \\
    \textbf{Dreieck-Stern-Umwandlung}
    \begin{align*}
        \underline{Z}_1 &= \frac{\underline{Z}_{12} \cdot \underline{Z}_{31}}{\underline{Z}_{12} + \underline{Z}_{23} + \underline{Z}_{31}} \\
        \underline{Z}_2 &= \frac{\underline{Z}_{12} \cdot \underline{Z}_{23}}{\underline{Z}_{12} + \underline{Z}_{23} + \underline{Z}_{31}} \\
        \underline{Z}_3 &= \frac{\underline{Z}_{23} \cdot \underline{Z}_{31}}{\underline{Z}_{12} + \underline{Z}_{23} + \underline{Z}_{31}} \\
        \text{für symm. Bel: } Z_Y &= \frac{Z_\Delta}{3} &
    \end{align*}

\subsubsection{Leistungen in DS}

    \begin{flalign*}
        \varphi &= \varphi_U - \varphi_I &
    \end{flalign*}

    \underline{Scheinleistung S [VA]:}
    \begin{flalign*}
        S &= 3 \cdot U_{LE} \cdot I_L = \sqrt{3} \cdot U_{LL} \cdot I_L = \sqrt{P^2 + Q^2} &\\
        \underline{S} &= 3 \cdot \underline{U}_{LE} \cdot \underline{I}_L^* = \sqrt{3} \cdot \underline{U}_{LL} \cdot \underline{I}_L^* = P+jQ &
    \end{flalign*}

    \textit{in Sternschaltung:}
    \begin{flalign*}
        \underline{S}_{3\sim} &= \frac{U_{LL}^2}{\underline{Z}_{LN}^*} \qquad
        \underline{S}_{1\sim} = \frac{U_{LL}^2}{3 \cdot \underline{Z}_{LN}^*} &
    \end{flalign*}
    \(\underline{S}_{3 \sim}:\) 3-Phasen-Scheinleistung\\

    \underline{Wirkleistung P [W]: Realteil}
    \begin{flalign*}
        P   &= S \cdot \cos \varphi
    \end{flalign*}

    \underline{Blindleistung Q [var]: Imaginärteil}
    \begin{flalign*}
        Q &= S \cdot \sin \varphi = P \cdot \tan \varphi
    \end{flalign*}
    \(Q > 0\) : induktiv \qquad \(Q < 0\) : kapazitiv

\subsubsection{Momentanleistungen in DS (einphasig)}
{\small
    \begin{flalign*}
        &p(t) = u(t) \cdot i(t) = \hat{u} \cos(\omega t + \varphi_u) \cdot \hat{i} \cos(\omega t + \varphi_i) &\\
             &= \hat{u} \cdot \hat{i} \left[\cos(\omega t + \varphi_u) \cdot \cos(\omega t + \varphi_i)\right] &\\
             &= U \cdot I \cdot \left[ \cos(2 \omega t + \varphi_u + \varphi_i) + \cos(\varphi_u - \varphi_i)\right] &\\
             &= U \cdot I \cdot \cos(\varphi) \left[1 + \cos(2(\omega t + \varphi_i))\right] - U \cdot I \cdot \sin(\varphi) \sin(2(\omega t + \varphi_i)) &\\
             &= P \cdot \left[1 + \cos(2(\omega t + \varphi_i))\right] - Q \cdot \sin(2(\omega t + \varphi_i)) &
    \end{flalign*}
}


\subsection{Klassifizierung und Struktur}
    \subsubsection{Spannungsbezeichnungen}

        \textcolor{red}{\textbf{Alle Spannungswerte sind als Leiter-Leiter-Spannungen angegeben, da dies die höchste auftretende ist.}}

        \begin{itemize}
            \item \textbf{Netznennspannung \(U_n\): } \\
                dient zur Kennzeichnung, meist gleich der Betriebsspannung
            
            \item \textbf{Betriebsspannung \(U_b\): } \\
                tatsächliche momentane Spannung, abhängig von der Belastung

            \item \textbf{Höchste Spannung für Betriebsmittel \(U_m\): } \\
                maximale Spannung, die ein Betriebsmittel aushalten muss

            \item \textbf{Bemessungsspannung eines Betriebsmittel \(U_r\): } \\
                Spannung, für die ein Betriebsmittel ausgelegt ist 
                \begin{flalign*}
                    U_r \ge U_m
                \end{flalign*}
        \end{itemize}